% =============================================================
%  ch13_sim_real_trennung.tex  --  Wo sich Simulation und Realität trennen
% =============================================================

\chapter{Wo sich Simulation und reale Implementierung trennen}

Das ist die architektonisch wichtigste Einsicht des ganzen Dokuments. Der Trick
von \textbf{ros2\_control}: deine High-Level-Knoten (Trajektorienplaner, Vision,
Kalman, Nav2) bleiben \emph{identisch}. Was sich ändert, ist nur die unterste
Schicht -- das \textbf{Hardware-Interface-Plugin}.

\begin{center}
\begin{tikzpicture}[node distance=8mm]
\node[node,fill=accent!8] (high)
  {High-Level-Knoten\\\scriptsize Trajektorie $\cdot$ Vision $\cdot$ Kalman $\cdot$ Nav2 $\cdot$ MoveIt2};
\node[node,below=of high] (ctrl)
  {controller\_manager\\\scriptsize joint\_trajectory\_controller, joint\_state\_broadcaster};
\node[node,below left=10mm and -10mm of ctrl,fill=accent!18] (sim)
  {Sim-Hardware-Interface\\\scriptsize gz\_ros2\_control};
\node[node,below right=10mm and -10mm of ctrl,fill=accent2!12,draw=accent2] (real)
  {Reales Hardware-Interface\\\scriptsize micro-ROS / Serial $\rightarrow$ ESP32 $\rightarrow$ TMC2209};
\node[node,below=8mm of sim] (gz)
  {Gazebo Harmonic\\\scriptsize Physik + Sensoren};
\node[node,below=8mm of real,draw=accent2,fill=accent2!8] (hw)
  {Realer Roboter\\\scriptsize Motoren, Encoder, Kamera};
\draw[flow] (high)--(ctrl);
\draw[flow] (ctrl)--(sim);
\draw[flow] (ctrl)--(real);
\draw[flow] (sim)--(gz);
\draw[flow] (real)--(hw);
\node[font=\scriptsize\itshape,accent] at ($(high)+(6.4,0)$)
  {\;identisch in Sim \& Real};
\end{tikzpicture}
\captionof{figure}{Sim-zu-Real: nur die unterste Schicht (Hardware-Interface)
wird getauscht. Alles darüber bleibt gleich.}
\end{center}

\noindent Konkret: In der Simulation lädt die URDF das Plugin
\code{gz\_ros2\_control/GazeboSimSystem}. Für die echte Maschine schreibst du
ein eigenes \code{hardware\_interface::SystemInterface}, das Sollwerte an den
ESP32 schickt (micro-ROS oder serielle Brücke) und Gelenkzustände
(Encoder/StallGuard) zurückliest. \textbf{Dieselben Controller, dieselben
Topics, dieselbe Trajektorie} -- du tauschst eine Plugin-Zeile und die
Hardware. Genauso bei der Kamera: simulierte Bildquelle vs.\ \code{v4l2\_camera}
-- die nachgelagerte Vision-Pipeline merkt keinen Unterschied.
